% Generated from locale/tr/content/first-order-logic/axiomatic-deduction/identity.tex; do not hand-edit.


\olfileid[tr]{fol}{axd}{ide}

\olsection{\usetoken{P}{derivation} ve \usetoken{S}{identity}}

$\eq$ simgesini kullanan !!{derivation}s elde etmek için yalnızca yeni aksiyom
şemaları ekleriz. !!{derivation} ile $\Proves$ tanımları değişmeden kalır;
sadece yeni aksiyomlara da izin veririz.

\begin{defn}[!!{identity} aksiyomları]
Kapalı terimler $t$, $t_1$ ve~$t_2$ için
\begin{align}
\ollabel{ax:id1} & \eq[t][t], \\
\ollabel{ax:id2} & \eq[t_1][t_2] \lif (!B(t_1) \lif !B(t_2)).
\end{align}
\end{defn}

\begin{prop}\ollabel{prop:sound}
  \olref{ax:id1} ve \olref{ax:id2} aksiyomları geçerlidir.
\end{prop}

\begin{proof}
  Alıştırma.
\end{proof}

\begin{prob}
\olref[fol][axd][ide]{prop:sound} önermesini ispatlayın.
\end{prob}

\begin{prop}\ollabel{prop:iden1}
 Her kapalı terim~$t$ ve her $\Gamma$ kümesi için $\Gamma \Proves \eq[t][t]$.
\end{prop}

\begin{prop}\ollabel{prop:iden2}
  Kapalı terimler $t_1$ ve~$t_2$ için, $\Gamma \Proves !A(t_1)$ ve
  $\Gamma \Proves \eq[t_1][t_2]$ ise
  $\Gamma \Proves !A(t_2)$.
\end{prop}

\begin{proof}
Şu !!{formula}, \olref{ax:id2}'nin bir örneğidir:
\[
(\eq[t_1][t_2] \lif (!A(t_1) \lif !A(t_2))).
\]
Sonuç \MP{}'nin iki kez uygulanmasıyla elde edilir.
\end{proof}

